\subsection{Extension of Scalars in Semisimple Modules}

PROPOSITION 8. --- Let A be a K-algebra, M an A-module, and L an extension of the field K. Denote the commutant of M by D and the center of D by Z.

\begin{enumerate}
\def\labelenumi{\alph{enumi})}
\item
  Suppose that the \(A_{(L)}\) -module \(M_{(L)}\) is simple (resp. isotypical, semisimple). Then the A-module M is simple (resp. isotypical, semisimple).
\item
  Suppose that the A-module M is semisimple and that M or L is finite-dimensional over K. The \(\mathrm{A}_{(\mathrm{L})}\) -module \(\mathrm{M}_{(\mathrm{L})}\) is semisimple if and only if the ring \(\mathrm{Z}_{(\mathrm{L})}\) is reduced. The \(\mathrm{A}_{(\mathrm{L})}\) -module \(\mathrm{M}_{(\mathrm{L})}\) is isotypical and nonzero if and only if the ring \(\mathrm{Z}_{(\mathrm{L})}\) is an integral domain.
\item
  Suppose that the A-module M is simple. The \(\mathrm{A}_{(\mathrm{L})}\) -module \(\mathrm{M}_{(\mathrm{L})}\) is semisimple (resp. isotypical, simple) if and only if the ring \(\mathrm{D}_{(\mathrm{L})}\) is semisimple (resp. simple, a field).
\end{enumerate}

Assertion a) is a specific case of Proposition 1 (VIII, p.~211), assertion b) is a specific case of Proposition 5 (VIII, p.~219), and assertion c) is a specific case of Corollary 1 of VIII, p.~215.

COROLLARY 1. --- a) Suppose that the A-module M is semisimple, that the extension L of K is separable, and that M or L is finite-dimensional over K. Then the \(\mathrm{A}_{(\mathrm{L})}\) -module \(\mathrm{M}_{(\mathrm{L})}\) is semisimple.

\begin{enumerate}
\def\labelenumi{\alph{enumi})}
\setcounter{enumi}{1}
\tightlist
\item
  Suppose that the A-module M is simple and that its commutant is equal to K. Then the \(\mathrm{A}_{(\mathrm{L})}\) -module \(\mathrm{M}_{(\mathrm{L})}\) is simple.
\end{enumerate}

Assertion a) follows from Corollary VIII, p.~221. Assertion b) is a specific case of Proposition 8, c).

COROLLARY 2. --- Let L be an extension of the field K. Denote the center of the K-algebra A by Z.

\begin{enumerate}
\def\labelenumi{\alph{enumi})}
\item
  If the L-algebra \(A_{(L)}\) is semisimple, then the K-algebra A is semisimple.
\item
  Suppose that the K-algebra A is semisimple and that L or A is finite-dimensional over K. The L-algebra \(\mathrm{A}_{(\mathrm{L})}\) is semisimple if and only if the ring \(\mathrm{Z}_{(\mathrm{L})}\) is reduced; this is, in particular, the case when L is a separable extension of K. The L-algebra \(\mathrm{A}_{(\mathrm{L})}\) is simple if and only if the ring \(\mathrm{Z}_{(\mathrm{L})}\) is an integral domain; this is, in particular, the case when the center of A is equal to K.
\end{enumerate}

Assertions a) and b) follow from Proposition 8, a) and b) applied to the A-module \(\mathrm{A}_s\) .

PROPOSITION 9. --- Let A be a K-algebra and L a separable extension of K.

\begin{enumerate}
\def\labelenumi{\alph{enumi})}
\item
  If M is an A-module without radical, then the \(\mathrm{A}_{(\mathrm{L})}\) -module \(\mathrm{M}_{(\mathrm{L})}\) is without radical.
\item
  If the K-algebra A is without radical, then the L-algebra \(\mathrm{A}_{(\mathrm{L})}\) is without radical.
\end{enumerate}

Let us prove assertion a). Let M be an A-module without radical. We identify M with its canonical image in \(\mathrm{M}_{(\mathrm{L})}\) . Let N be a maximal submodule of M. Since the A-module M/N is simple, it follows from Proposition 4 of VIII, p.~218 that the \(\mathrm{A}_{(\mathrm{L})}\) -module \((\mathrm{M}/\mathrm{N})_{(\mathrm{L})} = \mathrm{M}_{(\mathrm{L})}/\mathrm{N}_{(\mathrm{L})}\) is without radical and therefore \(\mathfrak{R}_{\mathrm{A}_{(\mathrm{L})}}(\mathrm{M}_{(\mathrm{L})}) \subset \mathrm{N}_{(\mathrm{L})}\) by Corollary 1, c) of VIII, p.~152. Now, it follows from the corollary of Proposition 14 of II, §7, No.~7, p.~306 that the intersection of the \(\mathrm{N}_{(\mathrm{L})}\) , where N runs through the set of maximal submodules of M, is reduced to 0. Consequently, the \(\mathrm{A}_{(\mathrm{L})}\) -module \(\mathrm{M}_{(\mathrm{L})}\) is without radical.

Assertion b) follows from assertion a) applied to the A-module \(\mathrm{A}_s\) .

PROPOSITION 10. --- Let A be a K-algebra and L an extension of K. Let M be an A-module.

\begin{enumerate}
\def\labelenumi{\alph{enumi})}
\tightlist
\item
  We make one of the following two assumptions:
\end{enumerate}

\begin{enumerate}
\def\labelenumi{(\roman{enumi})}
\item
  The A-module M is finitely generated and L is algebraic over K.
\item
  The ring A is left Artinian.
\end{enumerate}

Then we have the inclusion

\[
\mathfrak {R} _ {\mathrm{A}} (\mathrm{M}) _ {(\mathrm{L})} \subset \mathfrak {R} _ {\mathrm{A} _ {(\mathrm{L})}} (\mathrm{M} _ {(\mathrm{L})}).
\]

\begin{enumerate}
\def\labelenumi{\alph{enumi})}
\setcounter{enumi}{1}
\tightlist
\item
  If L is a separable extension of K, then we have the inclusion
\end{enumerate}

\[
\mathfrak {R} _ {\mathrm{A} _ {(\mathrm{L})}} (\mathrm{M} _ {(\mathrm{L})}) \subset \mathfrak {R} _ {\mathrm{A}} (\mathrm{M}) _ {(\mathrm{L})}.
\]

Let us prove assertion a). Consider case (i). First suppose that L has finite degree over K. Then the A-module \(M_{(L)}\) is finitely generated. Denote the canonical homomorphism from A to \(A_{(L)}\) by f; the ring \(A_{(L)}\) is generated by the union of its center and \(f(A)\) . We can therefore apply Proposition 3 of VIII, p.~174 to the \(A_{(L)}\) -module \(M_{(L)}\) . This gives the inclusion \(\mathfrak{R}_{\mathrm{A}}(\mathrm{M}_{(\mathrm{L})}) \subset \mathfrak{R}_{\mathrm{A}_{(\mathrm{L})}}(\mathrm{M}_{(\mathrm{L})})\) and a fortiori \(\mathfrak{R}_{\mathrm{A}}(\mathrm{M})_{(\mathrm{L})} \subset \mathfrak{R}_{\mathrm{A}_{(\mathrm{L})}}(\mathrm{M}_{(\mathrm{L})})\) (VIII, p.~152, Corollary 1).

Let us now treat the general case. Let \(x_{1},\ldots,x_{n}\) be elements that generate the A-module M and x an element of \(\mathfrak{R}_{\mathrm{A}}(\mathrm{M})\) . Let \(a_{1},\ldots,a_{n}\) be elements of \(\mathrm{A}_{(\mathrm{L})}\) ; since L is algebraic over K, there exists a finite extension \(L'\) of K contained in L such that the \(a_{i}\) belong to \(\mathrm{A}_{(\mathrm{L}')}\) . By the above, x belongs to the radical of the \(\mathrm{A}_{(\mathrm{L}')}\) -module \(\mathrm{M}_{(\mathrm{L}')}\) ; it follows from the corollary of VIII, p.~153 that the elements \(x_{i}+a_{i}x\) for \(1\leqslant i\leqslant n\) generate the \(\mathrm{A}_{(\mathrm{L}')}\) -module \(\mathrm{M}_{(\mathrm{L}')}\) and therefore the \(\mathrm{A}_{(\mathrm{L})}\) -module \(\mathrm{M}_{(\mathrm{L})}\) . By the same corollary, x belongs to the radical of the \(\mathrm{A}_{(\mathrm{L})}\) -module \(\mathrm{M}_{(\mathrm{L})}\) .

Let us now consider case (ii). Let \(\mathfrak{r}\) be the radical of A, so that the radical of the A-module M is equal to \(\mathfrak{rM}\) (VIII, p.~174, Corollary). The radical \(\mathfrak{r}\) of A is a nilpotent two-sided ideal of A (VIII, p.~173, Proposition 1); therefore, \(\mathfrak{r}_{(L)}\) is a nilpotent two-sided ideal of \(A_{(L)}\) . It follows that \(\mathfrak{r}_{(L)} \subset \Re(A_{(L)})\) (VIII, p.~156, Theorem 1), and Proposition 6 of VIII, p.~158 implies \(\Re(A_{(L)})M_{(L)} \subset \Re_{A_{(L)}}(M_{(L)})\) . We therefore have \(\Re_A(M) = \mathfrak{rM} \subset \Re_{A_{(L)}}(M_{(L)})\) , which concludes the proof of assertion a).

The module \(\mathrm{M}/\mathfrak{R}_{\mathrm{A}}(\mathrm{M})\) is without radical. If L is a separable extension of K, then it follows from Proposition 9 of VIII, p.~223 that the \(\mathrm{A}_{(\mathrm{L})}\) -module \((\mathrm{M}/\mathfrak{R}_{\mathrm{A}}(\mathrm{M}))_{(\mathrm{L})}\) is without radical. Consequently, we have the inclusion

\[
\mathfrak {R} _ {\mathrm{A} _ {(\mathrm{L})}} (\mathrm{M} _ {(\mathrm{L})}) \subset \mathfrak {R} _ {\mathrm{A}} (\mathrm{M}) _ {(\mathrm{L})}.
\]

COROLLARY. --- Let L be a separable extension of K. We have \(\mathfrak{R}(A_{(L)}) = \mathfrak{R}(A)_{(L)}\) if L is algebraic over K or if the ring A is left Artinian.

This is the case \(\mathrm{M} = \mathrm{A}_s\) of Proposition 10.

